\documentclass[a4paper,12pt]{article}

\usepackage{ucs}
\usepackage[utf8x]{inputenc}
\usepackage[english,russian]{babel}
\usepackage{cmlgc}
\usepackage{graphicx}
\usepackage{geometry}
\usepackage{setspace}
\usepackage{fancyvrb}
\usepackage{appendix}
\usepackage{listings}
\usepackage{color}
\usepackage{amssymb}
\usepackage{hyphenat}
\makeatletter
\renewcommand\@biblabel[1]{#1.}
\makeatother

\newcommand{\myrule}[1]{\rule{#1}{0.4pt}}
\newcommand{\sign}[2][~]{{\small\myrule{#2}\\[-0.7em]\makebox[#2]{\it #1}}}

% Поля
\usepackage[top=20mm, left=30mm, right=15mm, bottom=20mm, includefoot]{geometry}
\usepackage{indentfirst}
\footskip 20pt

\definecolor{gray}{rgb}{0.5,0.5,0.5}
\setcounter{secnumdepth}{4}

\bibliographystyle{gost780s}  %% стилевой файл для оформления по ГОСТу

\lstset{
  basicstyle=\footnotesize,           % the size of the fonts that are used for the code
  numbers=left,                   % where to put the line-numbers
  numberstyle=\tiny\color{gray},
  captionpos=b,                   % sets the caption-position to bottom
  breaklines=true,                % sets automatic line breaking
}

\begin{document}

%%%%%%%%%%%%%%%%%%%%%%%%%%%%%%%
%%%                         %%%
%%% Начало титульного листа %%%

\thispagestyle{empty}
\begin{center}


\renewcommand{\baselinestretch}{1}
{\large
{\sc Петрозаводский государственный университет\\
Институт математики и информационных технологий\\
	Кафедра информатики и математического обеспечения
}
}

\end{center}


\begin{center}

 01.04.02 - Прикладная математика и информатика \\
	Интеллектуальные интернет-технологии
\end{center}

\vfill

\begin{center}
{\normalsize Отчет о практике по научно-исследовательской работе} \\
\medskip

%%% Название работы %%%
	Модели компьютерного зрения в производственных веб-системах \\
	(промежуточный)
\end{center}

\medskip

\begin{flushright}
\parbox{11cm}{%
\renewcommand{\baselinestretch}{1.2}
\normalsize
	Выполнил:\\
%%% ФИО студента %%%
студент 5 курса группы 22503
\begin{flushright}
	О. А. Плугин \sign[подпись]{4cm}
\end{flushright}
%%%%%%%%%%%%%%%%%%%%%%%%%
% девушкам применять "Выполнила" и "студентка"
%%%%%%%%%%%%%%%%%%%%%%%%%



Место прохождения практики: \\ 
Кафедра информатики и математического обеспечения\\

Период прохождения практики: \\
% Раскмментируйте строку с нужным Вам периодом прохождения.
%	\textcolor{red}{Для 4 к. бакалавриата и 6 к. магистратуры:} \\
% Период прохождения практики: \\ 02.09.19-15.12.19 \\
%	\textcolor{red}{Для остальных курсов:} \\
04.02.26 - 07.06.26 \\

% Если руководителей два - то раскомментровать строку про второго руководителя и применть "Руководители:"

Руководитель:\\
%%% степень, звание ФИО научного руководителя %%%
% Первый руководитель 
Д. Ж. Корзун, к. ф-м. н., доцент \\
\begin{flushright}
\sign[подпись]{4cm}
\end{flushright}


% Второй руководитель 
% \textcolor{red}{И. О. Фамилия, ученая степень, ученое звание} \\
% \begin{flushright}
% \sign[подпись]{4cm}
% \end{flushright}



Итоговая оценка
\begin{flushright}
  \sign[оценка]{4cm}
\end{flushright}
}
\end{flushright}

\vfill

\begin{center}
\large
    Петрозаводск --- 2026
\end{center}

%%% Конец титульного листа  %%%
%%%                         %%%
%%%%%%%%%%%%%%%%%%%%%%%%%%%%%%%

% Размер шрифта (14пт)
\large

% Межстрочный интервал (полуторный)
\onehalfspacing

\newpage

\tableofcontents

\newpage


%%%%%%%%%%%%%%%%%%%%%%%%%%%%%%%%
%%%                          %%%
%%% Введение                 %%%

%%% В введении Вы должны описать предметную область, с которой связана %%%
%%% Ваша работа, показать её актуальность, вкратце определить цель     %%%
%%% исследования/разработки					       %%%

\section*{Введение}
\addcontentsline{toc}{section}{Введение}

В условиях активного развития технологий и внедрения подходов искусственного интеллекта в прикладные задачи всё большую актуальность приобретают системы автоматизированного анализа визуальных данных. Компьютерное зрение находит широкое применение в промышленности, транспортной инфраструктуре и системах мониторинга, где требуется обработка больших объёмов видеоданных в режиме, близком к реальному времени. Особое значение при этом приобретают вопросы точности, производительности и адаптации алгоритмов машинного обучения к реальным условиям эксплуатации.

Одной из ключевых проблем при внедрении систем видеоаналитики является ограниченность вычислительных ресурсов, на которых требуется выполнять обработку видеоданных. Во многих прикладных задачах модели нейросетей должны функционировать на низкопроизводительных или мобильных устройствах, что накладывает ограничения на размер модели, время обработки одного кадра и общее потребление вычислительных ресурсов. В связи с этим актуальной является задача оптимизации моделей нейросетей компьютерного зрения с целью обеспечения возможности их применения в условиях ограниченных аппаратных ресурсов без существенной потери точности распознавания.

Второй важной практической задачей является автоматизация анализа видеоданных в предметных областях, где традиционно используется ручная или экспертная оценка. К таким задачам относится мониторинг состояния дорожной инфраструктуры, в частности автоматическое обнаружение и классификация дефектов дорожного покрытия по изображениям и видео. Использование методов компьютерного зрения в данной области позволяет повысить объективность оценки состояния дорог и сократить трудоёмкость обследований. Современные подходы к детекции объектов в режиме реального времени основаны на архитектурах YOLO, SSD и их модификациях \cite{redmon, liu}.


Таким образом, в рамках магистерской работы рассматриваются две взаимосвязанные практические задачи:
\begin{itemize}
    \item обеспечение возможности работы обученных моделей нейросетей компьютерного зрения на низкопроизводительных устройствах за счёт применения и анализа алгоритмов оптимизации;
    \item применение методов компьютерного зрения для автоматического распознавания дефектов дорожного покрытия по видеоданным.
\end{itemize}

Целью данной магистерской работы является анализ предметной области и современных подходов к применению и оптимизации методов компьютерного зрения для задач видеоаналитики, а также формирование теоретического и методического задела для дальнейших исследований в данной области.

Параллельно с теоретическим анализом ведётся работа над практическим проектом в области видеоаналитики дорожной инфраструктуры. Проект ориентирован на исследование подходов к автоматическому обнаружению и классификации дефектов дорожного покрытия на основе видеоданных, получаемых с мобильной лаборатории, оснащённой видеокамерами и дополнительными измерительными датчиками. Особое внимание при этом уделяется вопросам применимости и эффективности моделей нейросетей в реальных условиях эксплуатации.

Для достижения поставленной цели в рамках текущего этапа магистерской работы необходимо решить следующие взаимосвязанные задачи:
\begin{enumerate}
    \item Изучение и анализ научных подходов к оптимизации моделей нейросетей компьютерного зрения с точки зрения повышения производительности и снижения вычислительных затрат.
    \item Анализ практической задачи видеоаналитики дефектов дорожного покрытия, включая используемые данные, методы распознавания и требования к вычислительной эффективности моделей.
    \item Формирование выводов и определение направлений дальнейших исследований в рамках магистерской диссертации.
\end{enumerate}

Данный промежуточный отчёт имеет следующую структуру.

В \textbf{первом разделе} представлен анализ научной статьи, посвящённой методам оптимизации моделей нейросетей для задач компьютерного зрения.

Во \textbf{втором разделе} приводится описание проекта системы видеоаналитики дефектов дорожного покрытия и анализ применяемых подходов.

\newpage
\section{Методы повышения точности и производительности малых моделей детекции объектов}

\subsection{Постановка задачи}

В анализируемой научной статье рассматривается прикладная задача повышения эффективности модели нейросетей детекции объектов, предназначенной для обработки видеопотока в режиме реального времени на устройствах с ограниченными вычислительными ресурсами.

Основная практическая проблема заключается в том, что современные модели распознавания объектов обладают высокой точностью, но требуют значительных вычислительных ресурсов, что делает их непригодными для использования на мобильных и встраиваемых устройствах. В то же время упрощение архитектуры модели приводит к заметному снижению качества распознавания.

В качестве исходных ограничений рассматриваются процессор с 4 ядрами частотой 2,2 ГГц, оперативная память до 4 ГБ и обработка изображений с разрешением до 1920×1080 пикселей. При таких условиях использование крупных моделей нейросетей становится невозможным. Это показывает необходимость адаптации модели под реальные условия эксплуатации, и изменения включают в себя:
\begin{itemize}
    \item уменьшение числа слоёв и их глубины;
    \item изменения параметров входных параметров;
\end{itemize}

Таким образом, задача статьи формулируется как поиск баланса между точностью распознавания и производительностью модели, позволяющего использовать алгоритмы компьютерного зрения в реальных системах с ограниченными вычислительными ресурсами ресурсами.

В рамках данной главы методы повышения точности и производительности малых моделей детекции объектов рассматриваются на примере научной статьи \cite{taema}, в которой предложен алгоритм оптимизации параметров модели нейросетей для работы в условиях ограниченных вычислительных ресурсов.


\subsection{Параметры и метрики задачи}

В задаче оптимизации модели нейросетей рассматриваются следующие ключевые параметры:

\begin{itemize}
    \item \textbf{Архитектура сети:} количество слоёв $n$, шаг карты признаков $feat_{st}$ и количество фильтров $f$. Эти параметры влияют на точность распознавания $mAP$ и скорость работы модели $FPS$.
    \item \textbf{Входные данные:} размер входного изображения $W \times H$ ($W = H$). Любое изображение автоматически преобразуется к заданному размеру. Увеличение разрешения повышает точность, но снижает производительность.
    \item \textbf{Якори (anchors):} количество рамок $k$ и их размеры $w_a$, $h_a$, которые определяют качество локализации объектов (IoU).
\end{itemize}

\textbf{Метрики оценки}

Основные показатели качества и производительности модели:  

\begin{itemize}
    \item $mAP$ (mean Average Precision, \% \ 0-1) — среднее значение точности по всем классам и порогам IoU;
    \item $Precision$ (доля корректно найденных объектов, \% / 0–1);
    \item $Recall$ (доля реальных объектов, найденных моделью, \% / 0–1);
    \item $IoU$ (Intersection over Union, 0–1) — точность совпадения предсказанной области с истинной;
    \item $FPS$ (frames per second, кадры/сек) — количество обрабатываемых кадров в секунду;
    \item $T_1$ (мс) — время обработки одного кадра.
\end{itemize}

\(\frac{кадры}{сек}\)

\subsection{Формулировка задачи оптимизации}

Данную задачу можно представить как задачу оптимизации. Основные параметры модели:

\begin{itemize}
    \item $W, H$ — ширина и высота входного изображения (пиксели);
    \item $n$ — количество слоёв нейронной сети;
    \item $k$ — количество якорных рамок;
    \item $mAP_{raw}$ — нижняя граница mAP для исходного набора параметров $(W,H,k)$;
    \item $\Delta_{raw}$ — допустимое снижение точности;
    \item $\Delta$ — желаемый прирост точности (в процентах);
    \item $FPS_{min}$ — минимальное значение FPS для задачи;
    \item $FPS'$ — нормализованное FPS в диапазоне $[0,1]$, $FPS' = 1 - \frac{FPS}{\infty}$;
    \item $\alpha_1, 1-\alpha_1$ — коэффициенты значимости метрик.
\end{itemize}

Критерий оптимизации при заданном наборе параметров $(W,H,n,k)$ имеет вид:

\begin{equation} \label{opt_crit}
| 1 - (\alpha_1 \cdot mAP + (1 - \alpha_1) \cdot FPS') | \to \min,
\end{equation}
\[
\quad \alpha_1, mAP, FPS' \in [0,1].
\]

Ограничения для задачи:

\begin{equation} \label{n_rest}
n < n_{max},
\end{equation}
\begin{equation} \label{map_raw_rest}
mAP_{raw} = mAP_{n_0} - \Delta_{raw},
\end{equation}
\begin{equation} \label{map_rest}
mAP(W,H,n,k) \ge mAP_{raw} + \Delta,
\end{equation}
\begin{equation} \label{fps_rest}
FPS(W,H,n,k) \ge FPS_{min}.
\end{equation}

\subsection{Алгоритм повышения точности сети}

В работе \cite{taema} предложен итеративный алгоритм подбора архитектурных и входных параметров модели нейросетей, направленный на достижение компромисса между точностью детекции объектов и производительностью системы.


\subsubsection*{Изменение архитектуры}
\begin{enumerate}
    \item Зафиксировать начальное количество слоёв: $n_0$.
    \item Уменьшить количество слоёв до $n < n_0$ \ref{n_rest} методом модифицированного двоичного поиска:
    \begin{enumerate}
        \item Установить диапазон поиска: $n \in [n_{min}; n_{max}]$, где $n_{min} = 5$, $n_{max} = n_0 - 1$.
        \item Вычислить среднее значение: $n_{mid} = (n_{min} + n_{max})/2$.
        \item Измерить метрики: $mAP(n_{mid})$, $FPS(n_{mid})$.
        \item Если $FPS(n_{mid}) \ge FPS_{min}$ \ref{fps_rest} и $mAP(n_{mid}) \ref{map_raw_rest} \ge mAP_{raw}$, уменьшаем глубину модели: $n_{max} = n_{mid}$.
        \item Иначе $n = n_{max}$ (выбираем последнее удовлетворяющее значение).
    \end{enumerate}
\end{enumerate}

\subsubsection*{Изменение входных данных}
\begin{itemize}
    \item Размер входного изображения $W_n \times H_n$ изменяется кратно 32 (например, уменьшение: $640 \rightarrow 416 \rightarrow 320$, увеличение: $416 \rightarrow 512 \rightarrow 640$).
    \item Для повышения точности $mAP\uparrow$ и снижения производительности $FPS\downarrow$ увеличивают $W$, $H$.
    \item Для снижения точности $mAP\downarrow$ и повышения производительности $FPS\uparrow$ уменьшают $W$, $H$.
\end{itemize}

\textbf{Пример изменения размера входного изображения и соответствующей точности:}

\begin{center}
\begin{tabular}{c c}
\hline
Размер $W \times H$ & $mAP$ \\
\hline
$416 \times 416$ & 0.37 \\
$608 \times 608$ &  0.58 \\
\hline
\end{tabular}
\end{center}

\subsubsection*{Изменение количества якорных рамок}
Нахождение оптимального количества якорных рамок включает несколько этапов:

\begin{enumerate}
    \item \textbf{Масштабирование рамок:}  
    Каждая рамка исходного размера $(w_t, h_t)$ масштабируется под новое разрешение входного изображения $(W_i, H_i)$:
    \begin{equation}
    (w_r, h_r) = \left( \frac{W_n}{W_i} \cdot w_t, \frac{H_n}{H_i} \cdot h_t \right)
    \end{equation}
    где $(W_n, H_n)$ — исходное разрешение сети.
    
    \item \textbf{Кластеризация $k$-средних:}  
    Провести кластеризацию $(w_r, h_r)$ для $k$ центров:
    \begin{equation}
    (w_c, h_c) = k\text{-means}(w_r, h_r)
    \end{equation}
    где $(w_c, h_c)$ — усреднённые размеры объектов в кластерах.
    
    \item \textbf{Определение оптимального числа кластеров $k$:}  
    Экспериментируют со значением $k$ от 3 до 15, вычисляя силуэт средних коэффициентов:
    \begin{equation}
    \bar{S}(k) = \frac{1}{m} \sum_{i=1}^{m} S(i), 
    \end{equation}
    \begin{equation}
    \quad S_k = \frac{b - a}{\max(a,b)} \tag{8,9}
    \end{equation}
    где $m$ — число элементов в кластеризации, $a$ — среднее внутрикластерное расстояние, $b$ — ближайшее межкластерное расстояние.  
    Максимум $\bar{S}(k)$ определяет оптимальное количество кластеров $k$.
\end{enumerate}

\subsubsection*{Нормализация якорей}
Для корректной работы сети размеры якорных рамок нормализуются относительно шага карты признаков:
\begin{equation}
(w_a, h_a) = \left( \frac{w_c}{feat_{st}}, \frac{h_c}{feat_{st}} \right)
\end{equation}

где $w_a, h_a$ — размеры якорей в масштабе feature map, с которым работает последний слой модели,  
а $feat_{st} = 2^n$ — шаг карты признаков, где сеть содержит $n$ слоёв max pooling.  
Данный подход позволяет адаптировать размеры якорных рамок к распределению объектов в обучающем наборе, повышая точность локализации (IoU) без увеличения вычислительной нагрузки.

\subsubsection*{Оценка результатов}
После настройки архитектуры, размеров входного изображения и якорей проводится измерение основных метрик модели: $mAP$, $Precision$, $Recall$, $IoU$ и $FPS$.  
Эти показатели позволяют оценить баланс между точностью распознавания и производительностью модели, а также подтвердить эффективность проведённой оптимизации.

\subsubsection*{Порядок выполнения алгоритма}
Алгоритм повышения точности сети выполняется в следующем порядке:

\begin{enumerate}
    \item Нахождение значения $n$, удовлетворяющего условиям задачи.
    \item С найденным $n$ осуществляется поиск подходящего размера входного изображения $W$, $H$.
    \item С изменёнными параметрами $W$, $H$, $n$ производится вычисление оптимального количества якорных рамок $k$.
\end{enumerate}

\subsection{Итоги анализа}
Представленный подход позволяет адаптировать архитектурные и входные параметры модели нейросетей под доступные вычислительные ресурсы, обеспечивая:  

\begin{itemize}
    \item \textbf{Сопоставимые показатели:} достижение значений метрик качества, сопоставимых с исходной моделью, при оптимизации архитектуры.
    \item \textbf{Эффективность ресурсов:} значительное снижение вычислительной сложности модели, что важно для систем компьютерного зрения на мобильных и малопроизводительных устройствах.
\end{itemize}

\subsubsection*{Пример результатов работы алгоритма}

\begin{center}
\normalsize
\setlength{\tabcolsep}{4pt} % уменьшение ширины колонок
\begin{tabular}{|l|c|c|c|c|c|c|c|c|}
\hline
Модель & $n$ & $k$ & $W \times H$ & Recall (\%) & Precision (\%) & IoU (\%) & mAP (\%) & FPS \\
\hline
Базовая & 24 & 3 & 416×416 & 80 & 78 & 75 & 72 & 5–20 \\
Оптим. & 9 & 9 & 608×608 & 85 & 84 & 82 & 82.21 & 50–100 \\
\hline
\end{tabular}
\end{center} \cite{taema}

\textbf{Комментарий к результатам:}  

\begin{itemize}
    \item Число слоёв $n$ уменьшено с 24 до 9, что существенно снизило вычислительную сложность модели.
    \item Метрики точности (Recall, Precision, IoU, mAP) выросли на 5–10\% по сравнению с исходной моделью.
    \item FPS увеличился в 5–10 раз, что демонстрирует высокую эффективность оптимизации для мобильных и малопроизводительных устройств.
\end{itemize}

\newpage

\section{Разработка системы видеоаналитики дефектов дорожного покрытия}

Данный раздел посвящён практическому описанию проекта системы автоматического обнаружения и классификации дефектов дорожного покрытия на основе видеоданных. В рамках текущего этапа основная деятельность была сосредоточена на реализации технической части грантового проекта, что потребовало плотной инженерной работы по созданию конвейера обработки данных и обучению нейросетевых моделей. Основная цель разрабатываемой системы заключается в обеспечении автоматизированного мониторинга состояния дорог для оперативного выявления повреждений и поддержки принятия решений по обслуживанию инфраструктуры. 

Для формализации описания практической реализации системы видеоаналитики дефектов дорожного покрытия необходимо определить ключевые понятия, составляющие основу методологии разработки современных систем компьютерного зрения на базе глубокого обучения. Данный ликбез фиксирует базовый инженерный и математический стек понятий, используемый в рамках описываемого проекта.

\begin{enumerate}
    \item \textbf{Выборка данных (Датасет / Dataset)} --- упорядоченная и репрезентативная совокупность информационных объектов (в данном контексте --- цифровых изображений или отдельных видеокадров), предназначенная для обучения, настройки и валидации математических моделей. В рамках сквозного процесса проектирования итоговый массив данных разделяется на три функциональных подмножества:
    \begin{itemize}
        \item \textit{Обучающая (тренировочная) выборка (Train dataset)} --- подмножество данных, на котором непосредственно осуществляется оптимизация внутренних параметров (весовых коэффициентов) нейронной сети посредством минимизации функции потерь;
        \item \textit{Валидационная выборка (Validation dataset)} --- независимый массив данных, применяемый в процессе обучения для мониторинга обобщающей способности модели, контроля переобучения (overfitting) и итеративного подбора оптимальных гиперпараметров;
        \item \textit{Тестовая выборка (Test dataset)} --- контрольное подмножество, полностью исключенное из процесса разработки и предназначенное для финальной объективной оценки метрик точности созданной системы в условиях, максимально приближенных к реальной эксплуатации.
    \end{itemize}

    \item \textbf{Конвейер обработки данных (Пайплайн / Data Pipeline)} --- структурированная и автоматизированная последовательность технологических этапов, обеспечивающая сквозное прохождение, трансформацию и диагностику информации от источника сырых данных (видеофайлов мобильной лаборатории) до вычислительного модуля нейросети. Типичный пайплайн в задачах видеоаналитики включает стадии декодирования видеопотока, фильтрации визуального шума, яркостной и геометрической нормализации, аннотирования, агрегации разрозненных источников, аугментации и финальной конвертации в тензорный формат, пригодный для параллельных вычислений на графическом процессоре (GPU).

    \item \textbf{Аннотирование (Разметка) данных} --- процесс сопоставления объектам на изображении метаданных, содержащих информацию об их пространственном положении (координатах) и семантическом классе. Для задач инстанс-сегментации (Instance Segmentation) разметка представляет собой детерминацию точных контуров (полигонов) каждого индивидуального дефекта, что позволяет модели не просто локализовать объект в прямоугольной рамке, но и вычислять его точную площадь и геометрию.

    \item \textbf{Инференс (Inference)} --- этап практической эксплуатации обученной нейросетевой модели, заключающийся в выполнении прямого прохода (forward pass) через сформированный вычислительный граф сети для генерации предсказаний (детекции и сегментации) на новых, ранее неизвестных системе видеоданных в режиме реального времени.
\end{enumerate}

\subsection{Цели и задачи системы}

Проектируемая система видеоаналитики направлена на решение комплекса научно-технических проблем, возникающих при потоковой обработке видеоданных в производственных условиях. Основные цели системы включают:
\begin{itemize}
    \item автоматическое обнаружение, инстанс-сегментацию и классификацию дефектов дорожного покрытия;
    \item формирование аналитической информации и объективной статистики о состоянии дорожного полотна;
    \item обеспечение возможности генерации регламентированных отчётов в форматах ГОСТ и ОДМ.
\end{itemize}

Для достижения поставленных целей в рамках грантового проекта был определен и реализован следующий пул задач:
\begin{enumerate}
    \item Обработка и нормализация сырых видеоданных, получаемых с аппаратных комплексов мобильной лаборатории.
    \item Формирование, программная диагностика и аугментация обучающей выборки.
    \item Детекция и сегментация дефектов с использованием оптимизированных архитектур глубоких сверточных нейронных сетей.
    \item Автоматизация контроля качества обучения модели и визуализация метрик (функции потерь, матрицы ошибок).
    \item Подготовка обученной модели к деплою в реальную производственную веб-систему с учетом ограничений аппаратных ресурсов.
\end{enumerate}

\subsection{Формирование обучающей выборки и пайплайн данных}

В задачах компьютерного зрения, базирующихся на машинном обучении, ключевым фактором, определяющим итоговую обобщающую способность модели, является качество набора размеченных данных. Техническая сложность проекта заключалась в формировании репрезентативного датасета, так как нейросетевая модель самостоятельно выводит математические закономерности целевых объектов исключительно на основе предоставленных примеров.

Основой для формирования обучающей выборки послужили видеоматериалы с зафиксированными повреждениями дорожного полотна. Первичная ручная разметка (аннотирование) видеокадров осуществлялась с использованием специализированной платформы CVAT (Computer Vision Annotation Tool). На данном этапе целевые объекты выделялись строгими полигонами для обеспечения последующей инстанс-сегментации.

Для интеграции полученных аннотаций в процесс обучения был разработан автоматизированный конвейер обработки данных (пайплайн). В ходе программной диагностики выгрузок из CVAT был утвержден финальный перечень, состоящий из 6 классов конечных дефектов. Поскольку процесс разметки производился итеративно и результаты сохранялись в виде множества разрозненных файлов, был реализован специализированный скрипт для их слияния. В рамках данного процесса было произведено объединение 15 задач разметки, содержащих более 58 000 изображений, в единый унифицированный массив с разделением на тренировочную и валидационную выборки.

Для обеспечения полной совместимости с современными фреймворками глубокого обучения итоговый консолидированный датасет был конвертирован в стандарт COCO. Данный формат является оптимальным стандартом де-факто для задач инстанс-сегментации, поскольку позволяет структурированно хранить точные координаты полигонов каждого обнаруженного объекта на кадре.

\subsection{Архитектура нейронной сети и процесс обучения}

Для решения задачи локализации и точного выделения контуров повреждений дорожного полотна была выбрана современная одностадийная архитектура нейронной сети RTMDet-Ins (Real-time Models for Object Detection and Instance Segmentation) \cite{lyu_rtmdet}. Данный выбор обусловлен высокой эффективностью архитектуры, обеспечивающей качество распознавания на уровне state-of-the-art при сохранении высокой скорости обработки кадров (FPS). Высокая производительность и сбалансированность модели критически важны для последующего развертывания вычислительного модуля на устройствах с ограниченными аппаратными ресурсами и для стабильного функционирования системы в веб-среде.

Практическая реализация и проведение серии вычислительных экспериментов потребовали развертывания специализированного программного окружения на базе фреймворков глубокого обучения PyTorch и MMDetection. В рамках настройки среды была выполнена адаптация конфигурационных файлов под специфику сформированного набора данных. Инженерные изменения включали переопределение путей к директориям данных, корректировку количества целевых классов под утвержденный перечень из 6 дефектов, а также конфигурирование параметров аугментации изображений для повышения устойчивости модели к внешним факторам.

Обучение нейронной сети выполнялось с применением методологии трансферного обучения (transfer learning). В качестве базовой инициализации использовались весовые коэффициенты модели RTMDet-Ins, предварительно обученной на масштабном наборе данных общего назначения COCO. \cite{lin_coco} Поскольку предобученная модель уже обладает сформированными внутренними фильтрами для выделения базовых геометрических признаков, контрастных границ, перспективы и форм, её применение позволило исключить этап обучения с нуля. Процесс обучения фактически представлял собой дообучение (fine-tuning) и узкую специализацию внутренних фильтров под конкретные типы разрушений дорожного покрытия. Это позволило радикально сократить вычислительный цикл, ускорить сходимость модели и гарантировать высокую точность детекции даже на относительно небольшом объеме специфических дорожных данных.

Управление процессом адаптации внутренних весов модели осуществлялось через итерационный подбор оптимальных гиперпараметров. Настройка включала в себя варьирование размера партии (batch size), определяющего количество одновременно анализируемых кадров перед обновлением весов, и регулирование скорости обучения (learning rate), задающего величину математического шага корректировки коэффициентов. Оптимизация процесса сходимости проводилась с привлечением современных алгоритмов обновления весов и механизмов регуляризации (weight decay). Использование регуляризации накладывало штрафы на чрезмерную сложность модели, не позволяя ей фокусироваться на случайных визуальных артефактах и бликах света, что принуждало систему анализировать обобщенные признаки дефектов для реальных условий эксплуатации.

\subsection{Оценка качества модели и анализ результатов}

Завершение этапа оптимизации весовых коэффициентов требует объективной оценки обобщающей способности обученной нейросети на независимой валидационной выборке. В рамках разрабатываемой системы видеоаналитики качество детекции и инстанс-сегментации дефектов оценивалось с использованием строгих математических метрик, исключающих субъективное визуальное восприятие. 

В качестве основных критериев эффективности применялись следующие показатели:
\begin{itemize}
    \item \textbf{Уверенность предсказания (Confidence Score)} --- внутренний показатель активации фильтров модели, применяемый для минимизации визуального шума и отсечения ложноположительных срабатываний посредством установки программного порога.
    \item \textbf{Пространственная точность (Intersection over Union, IoU)} --- метрика физической точности сегментации, вычисляющая отношение площади пересечения предсказанного полигона дефекта с истинным контуром из разметки к их суммарной площади.
    \item \textbf{Усредненная средняя точность (Mean Average Precision, mAP)} --- комплексная интегральная метрика, балансирующая показатели точности (Precision) и полноты (Recall) по всем классам дефектов. Высокий уровень mAP свидетельствует о способности системы обнаруживать подавляющее большинство повреждений дорожного полотна при минимальном количестве ложных срабатываний на визуальные артефакты.
\end{itemize}

Для обеспечения систематического контроля качества обучения и объективной оценки работы модели был разработан специализированный инструментарий аналитики. Вместо ручного разбора журналов обучения (логов) были написаны Python-скрипты, автоматизирующие генерацию аналитических отчетов и визуализацию метрик. Данный подход обеспечил непрерывный мониторинг процесса оптимизации.

В частности, было реализовано автоматическое построение графиков функции потерь (loss curves) для тренировочной и валидационной выборок (см. рис. \ref{fig:loss_graph}). Анализ динамики функции потерь позволил своевременно отслеживать процесс сходимости архитектуры сети и предотвращать эффект переобучения (overfitting), при котором модель заучивает специфические тренировочные данные, теряя способность к обобщению.
    
\begin{figure}[htbp]
    \centering
    \includegraphics[width=0.8\textwidth]{images/loss_graph.jpg}
    \caption{Динамика функции потерь (Loss) в процессе обучения модели RTMDet-Ins}
    \label{fig:loss_graph}
\end{figure}

\begin{figure}[htbp]
    \centering
    \includegraphics[width=0.8\textwidth]{images/mAP_graph.jpg}
    \caption{Динамика интегральной метрики mAP в процессе обучения модели RTMDet-Ins}
    \label{fig:map_graph}
\end{figure}

Дополнительно был внедрен функционал построения матриц ошибок (confusion matrix), пример которой представлен на рис. \ref{fig:conf_matrix}. Матрица обеспечила наглядное представление вероятностного распределения предсказаний модели по 6 целевым классам дефектов. Анализ матрицы позволил выявить проблемные паттерны классификации: степень дифференциации визуально схожих типов повреждений, а также преобладающие ложноположительные (False Positive) и ложноотрицательные (False Negative) срабатывания. Данная аналитика послужила фактологической основой для последующей балансировки обучающей выборки на новых итерациях дообучения модели.

\begin{figure}[htbp]
    \centering
     \includegraphics[width=0.8\textwidth]{images/confusion_matrix.png}
    \caption{Матрица ошибок (Confusion Matrix) классификации дефектов дорожного покрытия на валидационной выборке}
    \label{fig:conf_matrix}
\end{figure}

\subsection{Оптимизация инференса и подготовка к интеграции}

Заключительным этапом текущего технологического цикла разработки системы видеоаналитики является подготовка обученной нейросетевой модели RTMDet-Ins к деплою и последующей интеграции в реальную производственную веб-систему мониторинга состояния автомобильных дорог. Специфика эксплуатации моделей глубокого обучения в продакшн-среде предъявляет жесткие требования к вычислительной эффективности, латентности (времени отклика) и объему резервируемой памяти, что делает нецелесообразным использование стандартных исследовательских конфигураций.

По умолчанию фреймворк PyTorch сохраняет весовые коэффициенты и структуру вычислительного графа в собственном сериализованном формате, который ориентирован на возобновления сессий обучения, но является избыточным и тяжеловесным для этапа эксплуатации. В связи с этим в рамках проекта был выполнен детальный анализ способов конвертации итоговых весов модели инстанс-сегментации в открытый стандарт ONNX (Open Neural Network Exchange). \cite{bai_onnx}

Перевод обученной архитектуры RTMDet-Ins \cite{lyu_rtmdet} в формат ONNX позволяет реализовать комплексную оптимизацию этапа инференса по нескольким направлениям: \cite{bai_onnx}
\begin{itemize}
    \item \textbf{Декуплинг от среды разработки:} Экспорт вычислительного графа в формат ONNX позволяет полностью отвязать процесс исполнения (инференса) от исходного фреймворка PyTorch и библиотеки MMDetection. Это устраняет необходимость развертывания громоздких зависимостей на целевом сервере и упрощает архитектуру production-окружения.
    \item \textbf{Минимизация времени обработки:} Применение специализированного движка ONNX Runtime для исполнения модели существенно ускоряет время прямого прохода (forward pass) нейронной сети при обработке одного кадра. Высокая скорость вычислений критически важна для обеспечения стабильного фреймрейта (FPS) при анализе потокового видео высокого разрешения.
    \item \textbf{Эффективное управление памятью:} Среда выполнения ONNX Runtime на низком уровне оптимизирует структуру вычислительного графа, выполняя слияние операторов (operator fusion) и оптимизируя распределение памяти. Данный подход минимизирует нагрузку на оперативную память (RAM) и видеопамять (VRAM) устройства \cite{microsoft_onnxruntime}, позволяя запускать систему в веб-среде и на ресурсоемких локальных узлах.
\end{itemize}

На текущем этапе была полностью завершена теоретическая проработка пайплайна деплоя и сформированы требования к программному интерфейсу. Практическая реализация конвертера весов и его интеграция в веб-приложение мониторинга дорог входят в число приоритетных задач на ближайший период разработки. Создание этого механизма обеспечит бесшовное внедрение обученной нейросети в состав автономного аналитического модуля.

\newpage

\section*{Заключение}
\addcontentsline{toc}{section}{Заключение}

В рамках выполнения текущего этапа магистерской диссертации была проведена комплексная научно-исследовательская и практическая инженерная работа, направленная на разработку и оптимизацию системы автоматизированного визуального контроля дефектов дорожного покрытия. Поставленные перед данным этапом цели были успешно достигнуты путем решения комплекса взаимосвязанных задач.

В первой части работы был проведен детальный теоретический анализ современных подходов к оптимизации моделей нейросетей компьютерного зрения для функционирования на устройствах с ограниченными вычислительными ресурсами. Были рассмотрены методы снижения вычислительной сложности алгоритмов при сохранении высокой точности распознавания, проанализировано влияние ключевых параметров модели (количества слоёв, разрешения входных данных, якорей), а также формализован алгоритм поиска баланса между производительностью и качеством детекции. 

Во второй части работы был успешно реализован инженерный этап грантового проекта по созданию программного комплекса видеоаналитики. Выполнен полный цикл подготовки данных: от программной диагностики первичных аннотаций в среде CVAT до формирования унифицированного датасета и его конвертации в стандарт COCO для задач инстанс-сегментации. На базе фреймворков PyTorch и MMDetection проведено развертывание вычислительной среды и выполнено обучение современной архитектуры RTMDet-Ins. Применение методологии трансферного обучения позволило существенно ускорить сходимость модели и повысить точность обнаружения дефектов.

Для обеспечения объективного контроля качества обучения были разработаны инструменты автоматизированной аналитики, генерирующие графики функции потерь и матрицы ошибок, что позволило выявить проблемные классы для последующей балансировки датасета. Также была проведена успешная теоретическая проработка конвейера экспорта весовых коэффициентов обученной нейросети в кросс-платформенный формат ONNX.

Таким образом, на текущем этапе достигнуты следующие результаты:
\begin{itemize}
    \item систематизированы знания о методах оптимизации производительности моделей нейросетей;
    \item разработан автоматизированный пайплайн обработки данных и сформирован датасет в формате COCO;
    \item осуществлено обучение модели инстанс-сегментации RTMDet-Ins с применением Transfer Learning;
    \item внедрены Python-скрипты для автоматизированной оценки математических метрик качества (mAP, функции потерь, Confusion Matrix);
    \item проанализированы методы деплоя модели в производственную среду с использованием ONNX Runtime.
\end{itemize}

Определены следующие направления дальнейшей работы в рамках магистерской диссертации:
\begin{itemize}
    \item практическая реализация конвейера конвертации итоговых весов модели в формат ONNX;
    \item интеграция оптимизированного вычислительного модуля (инференса) в разрабатываемую веб-систему мониторинга дорожной инфраструктуры;
    \item проведение нагрузочных тестов и полевых испытаний алгоритмов на новых видеоданных мобильной лаборатории.
\end{itemize}

Выполненный этап работы позволил перевести проект из стадии сбора данных и теоретического анализа в стадию практической реализации вычислительного модуля, сформировав надежный технологический фундамент для дальнейшего внедрения системы в эксплуатацию.

\newpage

\renewcommand{\bibname}{Список литературы}
\begin{thebibliography}{99}
%\thispagestyle{empty}
\vspace{5mm}
\addcontentsline{toc}{section}{Список литературы}

\bibitem{taema}
Tarek Teama, Hongbin Ma, Ali Maher, and Mohamed A. Kassab. 2019. Real Time Object Detection Based on Deep Neural Network. In Intelligent Robotics and Applications: 12th International Conference, ICIRA 2019, Shenyang, China, August 8–11, 2019, Proceedings, Part IV. Springer-Verlag, Berlin, Heidelberg, 493–504.

\bibitem{redmon}
Redmon, J. You Only Look Once: Unified, Real-Time Object Detection / J. Redmon, S. Divvala, R. Girshick, A. Farhadi. - Вашингтон: IEEE Conference on Computer Vision and Pattern Recognition CVPR, 2016.- С. 10.

\bibitem{liu}
SSD: Single Shot MultiBox Detector, Computer Vision – ECCV 2016 / W. Liu, D. Anguelov, D. Erhan [и др.]. - Амстердам: Springer International Publishing AG, 2016.- с. 21-37.

\bibitem{wang_maskrcnn}
Wang, P. Research on Automatic Pavement Crack Recognition Based on the Mask R-CNN Model / P. Wang, C. Wang, H. Liu, M. Liang, W. Zheng, H. Wang, S. Zhu, G. Zhong, S. Liu // Coatings. – 2023. – Vol. 13, No. 2. – Article 430. – DOI: 10.3390/coatings13020430.

\bibitem{tapamo_cnn}
Tapamo, H. Convolutional Neural Networks for Crack Detection on Flexible Road Pavements / H. Tapamo, A. Bosman, J. Maina, E. Horak. – arXiv preprint arXiv:2304.02933, 2023. – DOI: 10.48550/arXiv.2304.02933.

\bibitem{lyu_rtmdet}
Lyu, C. RTMDet: An Empirical Study of Designing Real-Time Object Detectors / C. Lyu, W. Zhang, H. Huang [и др.]. – arXiv preprint arXiv:2212.07784, 2022. – DOI: 10.48550/arXiv.2212.07784.

\bibitem{lin_coco}
Lin, T.-Y. Microsoft COCO: Common Objects in Context / T.-Y. Lin, M. Maire, S. Belongie [и др.]. – Computer Vision – ECCV 2014. – 2014. – С. 740-755.

\bibitem{bai_onnx}
Bai, J. ONNX: Open Neural Network Exchange / J. Bai, F. Lu, K. Zhang [и др.]. – GitHub repository. – 2019. – URL: https://github.com/onnx/onnx.

\bibitem{microsoft_onnxruntime}
Microsoft Research. ONNX Runtime: cross-platform, high-performance scoring engine for ML models / Microsoft Research. – Официальный технический релиз Microsoft. – URL: https://onnxruntime.ai/.

\end{thebibliography}

\end{document}
